\subsection{例：取整数值的多项式}

设 V 是有限维 \(\mathbf{Q}\)-向量空间，\(V^{*}\) 是其对偶，\(\mathcal{V}\) 是 V 中的一个格，\(\mathcal{V}^{*}\) 是对偶 \(\mathbf{Z}\)-模（其所对偶的格为 \(\mathcal{V}\)），它可典则地识别为 \(V^{*}\) 中的一个格，\(\mathbf{S}(V)\) 是 V 的对称代数，并且

\[
\lambda : \mathbf {S} (\mathrm{V}) \to \mathrm{A} (\mathrm{V} ^ {*})
\]

\(\mathbf{S}(\mathrm{V})\) 到 \(\mathrm{V}^*\) 上的多项式函数代数的典则双射（代数，第 IV 章，第 5 节，第 11 号，备注 1）。若将 \(\mathrm{A}(\mathrm{V}^* \times \mathrm{V}^*)\) 识别为 \(\mathrm{A}(\mathrm{V}^*) \otimes_{\mathbf{Q}} \mathrm{A}(\mathrm{V}^*)\)，则 \(\lambda\) 将 \(\mathbf{S}(\mathrm{V})\) 的余积变换为映射 \(\mathrm{A}(\mathrm{V}^*) \to \mathrm{A}(\mathrm{V}^* \times \mathrm{V}^*)\)，该映射把多项式函数 \(\varphi\)（定义在 \(\mathrm{V}^*\) 上）关联到多项式函数

\[
(x, y) \mapsto \varphi (x + y)
\]

在 \(\mathrm{V}^* \times \mathrm{V}^*\) 上（代数，第 IV 章，第 5 节，第 11 号，备注 2）。

记 \(\binom{\mathcal{V}}{\mathbf{Z}}\) 为 \(\mathbf{S}(\mathrm{V})\) 的子集，它由对应于从 \(\mathrm{V}^*\) 到 \(\mathbf{Q}\) 的多项式映射且在 \(\mathcal{V}^*\) 上取整数值的元素组成。

命题 2. (i) \(\begin{pmatrix} \mathcal{V} \\ \mathbf{Z} \end{pmatrix}\) 是双代数 \(\mathbf{S}(\mathrm{V})\) 中的双序，并且 \(\begin{pmatrix} \mathcal{V} \\ \mathbf{Z} \end{pmatrix} \cap \mathrm{V} = \mathcal{V}\)。

\begin{enumerate}
\def\labelenumi{(\roman{enumi})}
\setcounter{enumi}{1}
\item
  \(\mathbf{Z}\)-代数 \(\left( \begin{array}{c}\mathcal{V}\\ \mathbf{Z} \end{array} \right)\) 由 \(\left( \begin{array}{c}h\\ n \end{array} \right)\)（当 \(h\in \mathcal{V},n\in \mathbf{N}\) 时）生成。
\item
  若 \((h_1, \ldots, h_r)\) 是 \(\mathcal{V}\) 的一组基，则元素
\end{enumerate}

\[
\binom{h}{\mathbf {n}} = \binom{h _ {1}}{n _ {1}} \dots \binom{h _ {r}}{n _ {r}},
\]

其中 \(\mathbf{n} = (n_1, \ldots, n_r)\) 属于 \(\mathbf{N}^r\)，这些元素构成 \(\mathbf{Z}\)-模 \(\binom{\mathcal{V}}{\mathbf{Z}}\) 的一组基。

对 \(m \in \mathbf{N}\)，置 \(\mathbf{S}_{m}(\mathrm{V}) = \sum_{i \leq m} \mathbf{S}^{i}(\mathrm{V}), \mathbf{S}_{m}(\mathcal{V}) = \sum_{i \leq m} \mathbf{S}^{i}(\mathcal{V})\)。根据代数，第 IV 章，第 5 节，第 9 号，命题 15 及备注，

\[
\mathbf {S} _ {m} (\mathcal {V}) \subset \mathbf {S} _ {m} (\mathrm{V}) \cap \binom{\mathcal {V}}{\mathbf {Z}} \subset \frac {1}{m !} \mathbf {S} _ {m} (\mathcal {V})
\]

因此 \(\binom{\mathcal{V}}{\mathbf{Z}}\cap\mathrm{V}=\mathcal{V}\)。由于 \(\mathbf{S}_{m}(\mathcal{V})\) 是 \(\mathbf{S}_{m}(\mathrm{V})\) 中的一个格，\(\mathbf{S}_{m}(\mathrm{V})\cap\binom{\mathcal{V}}{\mathbf{Z}}\) 也是 \(\mathbf{S}_{m}(\mathrm{V})\) 中的一个格。另一方面，\(\mathbf{S}_{m}(\mathrm{V})\cap\binom{\mathcal{V}}{\mathbf{Z}}\) 是 \(\mathbf{S}_{m+1}(\mathrm{V})\cap\binom{\mathcal{V}}{\mathbf{Z}}\) 的直和因子（因为商无挠），因而有一个补空间，且该补空间是自由的 \(\mathbf{Z}\)-模。由此可知 \(\binom{\mathcal{V}}{\mathbf{Z}}\) 是自由的 \(\mathbf{Z}\)-模。显然，这是代数 \(\mathbf{S}(\mathrm{V})\) 中的一个幺序。设 \((u_{n})_{n\in\mathbf{N}}\) 是 \(\mathbf{Z}\)-模 \(\binom{\mathcal{V}}{\mathbf{Z}}\) 的一组基。它也是 \(\mathbf{Q}\)-模 \(\mathbf{S}(\mathrm{V})\) 的一组基，并且，对于所有

\[
\varphi \in \mathbf {S} (\mathrm{V} \times \mathrm{V}) = \mathbf {S} (\mathrm{V}) \otimes_ {\mathbf {Q}} \mathbf {S} (\mathrm{V}),
\]

存在唯一的序列 \((v_{n})\)，其元素属于 S(V)，使得 \(\varphi=\sum u_{n}\otimes v_{n}\)。如上，将 S(V) 识别为 A(V*)，并将 S(V) \(\otimes\) S(V) 识别为 A(V* × V*)。若 \(\varphi\in\binom{\mathcal{V}\times\mathcal{V}}{\mathbf{Z}}\)，则多项式函数 \(x\mapsto\varphi(x,y)\) 属于 \(\binom{\mathcal{V}}{\mathbf{Z}}\)，对所有 \(y\in\mathcal{V}^{*}\) 都成立。由此可知 \(v_{n}(y)\in\mathbf{Z}\) 对所有 n 及所有 \(y\in\mathcal{V}^{*}\) 都成立，换言之，\(v_{n}\in\binom{\mathcal{V}}{\mathbf{Z}}\)。这证明余积将 \(\binom{\mathcal{V}}{\mathbf{Z}}\) 映到 \(\binom{\mathcal{V}}{\mathbf{Z}}\otimes_{\mathbf{Z}}\binom{\mathcal{V}}{\mathbf{Z}}\)。若 \(h\in\mathcal{V}\) 且 \(n\in\mathbf{N}\)，则 \(\binom{h}{n}\) 将 \(u\in\mathcal{V}^{*}\) 映到整数 \(\binom{u(h)}{n}\)，所以 \(\binom{h}{n}\in\binom{\mathcal{V}}{\mathbf{Z}}\)。断言 (iii) 现在由将定理 1 应用于交换李代数 V 而得，(ii) 随之成立。

推论。令 \(\mathrm{X}\) 为一个不定元。多项式 \(\binom{\mathrm{X}}{n}\)，其中 \(n \in \mathbf{N}\)，构成以下 \(\mathbf{Z}\)-模的一组基：该模由 \(\mathrm{P} \in k[\mathrm{X}]\) 且满足 \(\mathrm{P}(\mathbf{Z}) \subset \mathbf{Z}\) 的多项式组成。

若 \(\mathrm{P}(\mathbf{Z}) \subset \mathbf{Z}\)，则拉格朗日插值公式（代数，第 IV 章，第 2 节，第 1 号，命题 6）表明 P 的系数属于 \(\mathbf{Q}\)；因此可以假设 \(k = \mathbf{Q}\)，并将命题 2 应用于 \(V = \mathbf{Q}\)、\(\mathcal{V} = \mathbf{Z}\)。

\subsection{一些公式}

在本号中，A 表示一个有单位元的结合代数。如果 \(x \in A\)，我们记 ad x 而不记 \(ad_{A}x\)。

引理 2. 若 \(x, y \in \mathbf{A}\) 且 \(n \in \mathbf{N}\)，

\[
\frac {(\operatorname{ad} x) ^ {n}}{n !} y = \sum_ {p + q = n} (- 1) ^ {q} \frac {x ^ {p}}{p !} y \frac {x ^ {q}}{q !} = \sum_ {p + q = n} (- 1) ^ {q} x ^ {(p)} y x ^ {(q)}.\tag{11}
\]

事实上，若以 \(L_{x}\) 和 \(R_{x}\) 分别表示从 A 到 A 的映射 \(z \mapsto xz\) 和 \(z \mapsto zx\)，由于 \(L_{x}\) 与 \(R_{x}\) 对易，我们有

\[
\frac {1}{n !} (\mathrm{ad} x) ^ {n} = \frac {1}{n !} (\mathrm{L} _ {x} - \mathrm{R} _ {x}) ^ {n} = \sum_ {p + q = n} (- 1) ^ {q} \frac {1}{p !} \mathrm{L} _ {x} ^ {p} \frac {1}{q !} \mathrm{R} _ {x} ^ {q}.
\]

引理 3. 设 \(x, h \in \mathbf{A}\) 且 \(\lambda \in k\) 满足 (ad \(h)x = \lambda x\)。对于所有 \(n \in \mathbf{N}\) 及所有 \(\mathrm{P} \in k[\mathrm{X}]\)，有

\[
\mathrm{P} (h) x ^ {(n)} = x ^ {(n)} \mathrm{P} (h + n \lambda).\tag{12}
\]

由于 \(\operatorname{ad} h\) 是 \(A\) 的导子，且 \((\operatorname{ad} h)x\) 与 \(x\) 对易，我们有

\[
(\operatorname{ad} h) x ^ {n} = n x ^ {n - 1} ((\operatorname{ad} h) x) = n \lambda x ^ {n},\tag{13}
\]

所以

\[
(\operatorname{ad} h) x ^ {(n)} = n \lambda x ^ {(n)}.
\]

因此，公式 (12) 由特殊情形

\[
\mathrm{P} (h) x = x \mathrm{P} (h + \lambda)\tag{14}
\]

将 \(x\) 换为 \(x^{(n)}\)、将 \(\lambda\) 换为 \(n\lambda\) 即可。只需在 \(\mathrm{P} = \mathrm{X}^m\) 时证明 (14)，对 \(m\) 归纳即可。当 \(m = 0, 1\) 时显然。若 (14) 对 \(\mathrm{P} = \mathrm{X}^m\) 成立，则

\[
h ^ {m + 1} x = h. h ^ {m} x = h x (h + \lambda) ^ {m} = x (h + \lambda) ^ {m + 1}
\]

这就证明了 (12)。

引理 4. 设 \(x, y, h \in \mathbf{A}\) 使得

\[
[ y, x ] = h, \quad [ h, x ] = 2 x, \quad [ h, y ] = - 2 y.\tag{15}
\]

\begin{enumerate}
\def\labelenumi{(\roman{enumi})}
\tightlist
\item
  对于 \(m, n \in \mathbf{N}\)，有
\end{enumerate}

\[
x ^ {(n)} y ^ {(m)} = \sum_ {p \geq 0} y ^ {(m - p)} \binom{m + n - p - 1 - h}{p} x ^ {(n - p)}.\tag{16}
\]

\begin{enumerate}
\def\labelenumi{(\roman{enumi})}
\setcounter{enumi}{1}
\tightlist
\item
  令 \(\mathbf{A}'\) 为一个 \(\mathbf{Z}\)-子代数，且包含于 \(\mathbf{A}\)，它由 \(x^{(m)}\) 和 \(y^{(m)}\)（\(m \in \mathbf{N}\)）生成。于是 \(\binom{h}{n} \in \mathbf{A}'\)，对所有 \(n \in \mathbf{N}\) 都成立。
\end{enumerate}

公式 (16) 可写成等价形式

\[
(\operatorname{ad} x ^ {(n)}) y ^ {(m)} = \sum_ {p \geq 1} y ^ {(m - p)} \binom {m + n - p - 1 - h} {p} x ^ {(n - p)}.\tag{$17_{m}$}
\]

\(m = 0\) 时这是显然的。我们对 \(m\) 归纳。由 \((17_{m})\) 得到

\[
\begin{array}{l} (m + 1) (\operatorname{ad} x ^ {(n)}) y ^ {(m + 1)} = (\operatorname{ad} x ^ {(n)}) y ^ {(m)}. y + y ^ {(m)}. (\operatorname{ad} x ^ {(n)}) y \\ = \sum_ {p \geq 1} y ^ {(m - p)} \binom {m + n - p - 1 - h} {p} x ^ {(n - p)} y + y ^ {(m)} (n - 1 - h) x ^ {(n - 1)} \end{array} \tag {18}
\]

（第 1 节，第 1 号，引理 1）。现在应用同一引理，再应用引理 3，得到

\[
\begin{array}{l} \binom {m + n - p - 1 - h} {p} x ^ {(n - p)} y \\ = \binom {m + n - p - 1 - h} {p} (y x ^ {(n - p)} + (n - p - 1 - h) x ^ {(n - p - 1)}) \\ = y \binom {m + n - p + 1 - h} {p} x ^ {(n - p)} \\ \qquad + \binom {m + n - p - 1 - h} {p} (n - p - 1 - h) x ^ {(n - p - 1)}. \end{array}
\]

将其代入 (18)，得到

\[
\begin{array}{l} (m + 1) (\operatorname{ad} x ^ {(n)}) y ^ {(m + 1)} \\ \qquad = \sum_ {p \geq 1} (m - p + 1) y ^ {(m - p + 1)} \binom {m + n - p + 1 - h} {p} x ^ {(n - p)} \\ \qquad + \sum_ {p \geq 1} y ^ {(m - p)} \binom {m + n - p - 1 - h} {p} (n - p - 1 - h) x ^ {(n - p - 1)} \\ \qquad + y ^ {(m)} (n - 1 - h) x ^ {(n - 1)} \\ \qquad = \sum_ {p \geq 1} (m - p + 1) y ^ {(m - p + 1)} \binom {m + n - p + 1 - h} {p} x ^ {(n - p)} \\ \qquad + \sum_ {p \geq 0} y ^ {(m - p)} \binom {m + n - p - 1 - h} {p} (n - p - 1 - h) x ^ {(n - p - 1)}. \end{array}
\]

在第二个和式中将 \(p\) 换为 \(p - 1\)，再重新组合各项，得到

\[
(m + 1) (\operatorname{ad} x ^ {(n)}) y ^ {(m + 1)} = \sum_ {p \geq 1} y ^ {(m - p + 1)} A _ {p} x ^ {(n - p)}\tag{19}
\]

其中

\[
A _ {p} = (m - p + 1) \binom{m + n - p + 1 - h}{p} + (n - p - h) \binom{m + n - p - h}{p - 1}.
\]

置 \(z = m + n - p - h\)，则也可写成

\[
\begin{array}{l} A _ {p} = \frac {1}{p !} (m - p + 1) (z + 1) z (z - 1) \dots (z - p + 2) \\ \qquad + \frac {1}{(p - 1) !} (z - m) z (z - 1) \dots (z - p + 2) \\ \qquad = \frac {1}{p !} z (z - 1) \dots (z - p + 2) [ (m - p + 1) (z + 1) + p (z - m) ] \\ \qquad = (m + 1) \binom {z} {p} = (m + 1) \binom {(m + 1) + n - p - 1 - h} {p}. \end{array}
\]

将其代入 (19)，得到 \((17_{m + 1})\)，因而 (i) 成立。

假设 \(\binom{h}{p} \in \mathrm{A}'\) 对 \(p < n\) 成立。那么，对于所有 \(\mathrm{P} \in \mathbf{Q}[\mathrm{T}]\)（次数 \(< n\) 且满足 \(\mathrm{P}(\mathbf{Z}) \subset \mathbf{Z}\)），有 \(\mathrm{P}(h) \in \mathrm{A}'\)（第 4 号，命题 2 的推论）。因此，考虑 (16) 中 \(m = n\) 的情形，

\[
\begin{array}{c} (- 1) ^ {n} \binom{h}{n} = \binom{n - 1 - h}{n} \\ = - x ^ {(n)} y ^ {(n)} + \sum_ {p = 0} ^ {n - 1} y ^ {(n - p)} \binom{2 n - p - 1 - h}{p} x ^ {(n - p)} \in \mathrm{A} ^ {\prime}; \end{array}
\]

故 (ii) 由对 \(n\) 的归纳得到。

\subsection{分裂约化李代数的包络代数中的双序}

设 \(\mathfrak{g}\) 是 \(\mathbf{Q}\) 上的约化李代数，\(\mathfrak{h}\) 是 \(\mathfrak{g}\) 的可分裂嘉当子代数，且 \(R = R(\mathfrak{g},\mathfrak{h})\)（第 2 节，第 1 号，备注 5）。

定义 1. 若 \(\mathcal{H}\) 是 \(\mathfrak{h}\) 中的一个格，则称其为相对于 \(\mathfrak{g}\) 的许可格，只要对所有 \(\alpha \in \mathrm{R}\) 都有 \(H_{\alpha} \in \mathcal{H}\) 及 \(\alpha(\mathcal{H}) \subset \mathbf{Z}\)。

备注。1) 令 B 为 R 的一组基。格 \(\mathcal{H}\) 位于 \(\mathfrak{h}\) 中。它是许可格，当且仅当 \(H_{\alpha} \in \mathcal{H}\) 且 \(\alpha(\mathcal{H}) \subset \mathbf{Z}\) 对所有 \(\alpha \in B\) 都成立。

\begin{enumerate}
\def\labelenumi{\arabic{enumi})}
\setcounter{enumi}{1}
\item
  令 \(\mathfrak{c}\) 为 \(\mathfrak{g}\) 的中心。则格 \(\mathcal{H}\) 位于 \(\mathfrak{h}\) 中，并且是许可格，当且仅当 \(\mathrm{Q}(\mathrm{R}^{\vee}) \subset \mathcal{H} \subset \mathrm{P}(\mathrm{R}^{\vee}) \oplus \mathfrak{c}\)。此时，格 \(\mathcal{H} \cap \mathcal{D}\mathfrak{g}\) 是嘉当子代数 \(\mathfrak{h} \cap \mathcal{D}\mathfrak{g}\) 中的 \(\mathcal{D}\mathfrak{g}\) 的许可格。可能存在许可格 \(\mathcal{H}\) 满足 \(\mathcal{H} \neq (\mathcal{H} \cap \mathcal{D}\mathfrak{g}) \oplus (\mathcal{H} \cap \mathfrak{c})\)（参见第 13 节，第 1 号.IX）。
\item
  若 \(\mathfrak{g}\) 是半单的，则 \(\mathfrak{h}\) 中的许可格正是其子群 \(\mathcal{H}\)，即 \(\mathfrak{h}\) 的子群，且满足 \(Q(R^{\vee})\subset \mathcal{H}\subset P(R^{\vee})\)。
\end{enumerate}

在本号余下部分，我们固定一个分裂约化李代数 \((\mathfrak{g},\mathfrak{h})\)、\(\mathrm{R}=\mathrm{R}(\mathfrak{g},\mathfrak{h})\) 的一组基 B，并且对每个 \(\alpha\in B\) 取一对 \((x_{\alpha},y_{\alpha})\)，满足

\[
y _ {\alpha} \in \mathfrak {g} ^ {- \alpha}, x _ {\alpha} \in \mathfrak {g} ^ {\alpha}, [ y _ {\alpha}, x _ {\alpha} ] = H _ {\alpha}.\tag{20}
\]

若记 \(n_{+}\)（分别为 \(n_{-}\)）为由 \(x_{\alpha}\)（分别为 \(y_{\alpha}\)）生成的 g 的子代数，则我们知道（第 3 节，第 3 号，命题 9 (iii)）

\[
\mathfrak {g} = \mathfrak {n} _ {-} \oplus \mathfrak {h} \oplus \mathfrak {n} _ {+}\tag{21}
\]

\[
\mathrm{U} (\mathfrak {g}) = \mathrm{U} (\mathfrak {n} _ {-}) \otimes_ {\mathbf {Q}} \mathrm{U} (\mathfrak {h}) \otimes_ {\mathbf {Q}} \mathrm{U} (\mathfrak {n} _ {+})\tag{22}
\]

（其中 \(\mathrm{U}(\mathfrak{g}),\ldots\) 分别是 \(\mathfrak{g},\ldots\) 的包络代数）。

记 \(\mathcal{U}_{+}\) 为一个 \(\mathbf{Z}\)-子代数，包含于 \(\mathrm{U}(\mathfrak{n}_{+})\)，并由 \(x_{\alpha}^{(n)}\)（其中 \(\alpha \in B\) 且 \(n \in \mathbf{N}\)）生成。令 W 为 R 的 Weyl 群，\(R_{+}\) 为相对于 B 的正根集合。

引理 5. (i) \(\mathcal{U}_{+}\) 是向量空间 \(\mathrm{U}(\mathfrak{n}_{+})\) 中的一个格。

\begin{enumerate}
\def\labelenumi{(\roman{enumi})}
\setcounter{enumi}{1}
\tightlist
\item
  对所有 \(\alpha \in \mathrm{B}\)，有 \(\mathcal{U}_{+} \cap \mathrm{U}(\mathfrak{g}^{\alpha}) = \bigoplus_{n \in \mathbf{N}} \mathbf{Z}x_{\alpha}^{(n)}\)。
\end{enumerate}

按定义，\(\mathcal{U}_{+}\) 作为 \(\mathbf{Z}\)-模由以下元素生成

\[
x _ {\varphi} ^ {(\mathbf {n})} = \prod_ {1 \leq i \leq r} x _ {\varphi (i)} ^ {(n (i))}
\]

其中 \(r \in \mathbf{N}\)、\(\varphi = (\varphi(i)) \in \mathrm{B}^{r}\)，且 \(\mathbf{n} = (n(i)) \in \mathbf{N}^{r}\)。给代数 \(\mathrm{U}(\mathfrak{n}_{+})\) 赋予 \(\mathrm{Q}(\mathrm{R})\) 型的分次，使每个 \(\mathfrak{g}^{\alpha} (\alpha \in \mathrm{R}_{+})\) 都是次数为 \(\alpha\) 的齐次空间。前述类型的单项式 \(x_{\varphi}^{(\mathbf{n})}\) 是次数为

\[
\sum_ {1 \leq i \leq r} n (i) \varphi (i) \in \mathrm{Q} (\mathrm{R}).
\]

具有给定次数 q 的这类单项式数目有限，并且它们在 Q 上生成 \(\mathrm{U}(\mathfrak{n}_{+})\) 的次数为 q 的齐次分量。这证明 (i)。

若 \(\alpha \in \mathrm{B}\)，则 \(\mathcal{U}_{+} \cap \mathrm{U}(\mathfrak{g}^{\alpha})\) 包含于次数为 \(\alpha\) 的倍数的齐次分量之和中；因此，根据上面的结论，\(\mathcal{U}_{+} \cap \mathrm{U}(\mathfrak{g}^{\alpha})\) 由 \(x_{\varphi}^{(\mathbf{n})}\) 生成，其中 \(\sum n(i)\varphi(i) \in \mathbf{N}\alpha\)。这迫使 \(\varphi(i) = \alpha\) 对所有 \(i\) 都成立（因为 B 是 R 的一组基），所以

\[
x _ {\varphi} ^ {(\mathbf {n})} = x _ {\alpha} ^ {(n (1))} \dots x _ {\alpha} ^ {(n (r))} = \frac {(n (1) + \cdots + n (r)) !}{n (1) ! \dots n (r) !} x _ {\alpha} ^ {(n (1) + \dots + n (r))}.
\]

因此，\(\mathcal{U}_{+} \cap \mathcal{U}(\mathfrak{g}^{\alpha}) \subset \bigoplus_{n} \mathbf{Z}x_{\alpha}^{(n)}\)，故 (ii) 成立。

证毕。

在本号余下部分，若 E 和 F 是 \(U(\mathfrak{g})\) 的 Z-子模，我们用 E.F 表示 \(U(\mathfrak{g})\) 的 Z-子模；它由乘积 ab 生成，其中 \(a \in E, b \in F\)。

命题 3. 设 \(\mathcal{H}\) 是 \(\mathfrak{h}\) 中的一个许可格。令 \(\mathcal{U}_{+},\mathcal{U}_{-},\mathcal{U}_{0}\) 为 \(\mathbf{Z}\)-子代数，包含于 \(\mathrm{U}(\mathfrak{g})\)，它们分别由元素 \(x_{\alpha}^{(n)}\)（\(\alpha \in \mathrm{B}, n \in \mathbf{N}\)）、\(y_{\alpha}^{(n)}\)（\(\alpha \in \mathrm{B}, n \in \mathbf{N}\)）、\(\binom{h}{n}\)（\(h \in \mathcal{H}, n \in \mathbf{N}\)）生成。令 \(\mathcal{U}\) 为 \(\mathbf{Z}\)-子代数 \(\mathrm{U}(\mathfrak{g})\)，由 \(\mathcal{U}_{+},\mathcal{U}_{-},\mathcal{U}_{0}\) 生成。

\begin{enumerate}
\def\labelenumi{(\roman{enumi})}
\item
  \(\mathcal{U}\) 是双代数 \(\mathrm{U}(\mathfrak{g})\) 中的双序。
\item
  有 \(\mathcal{U} = \mathcal{U}_{-}. \mathcal{U}_{0}. \mathcal{U}_{+}\)、\(\mathcal{U} \cap \mathfrak{h} = \mathcal{H}\)，并且对所有 \(\alpha \in B\)，
\end{enumerate}

\[
\mathcal {U} \cap \mathfrak {g} ^ {\alpha} = \mathbf {Z} x _ {\alpha}, \quad \mathcal {U} \cap \mathfrak {g} ^ {- \alpha} = \mathbf {Z} y _ {\alpha}.
\]

由引理 5 和命题 2，\(\mathcal{U}_{+},\mathcal{U}_{-},\mathcal{U}_{0}\) 分别是 \(\mathbf{Q}\)-代数 \(\mathrm{U}(\mathfrak{n}_{+}),\mathrm{U}(\mathfrak{n}_{-}),\mathrm{U}(\mathfrak{h})\) 中的序，并且

\[
\binom{\pm h + q}{p} \in \mathcal {U} _ {0} \quad \text { for } h \in \mathcal {H}, q \in \mathbf {Z}, p \in \mathbf {N}.\tag{23}
\]

置 \(\mathcal{L} = \mathcal{U}_{-}. \mathcal{U}_{0}. \mathcal{U}_{+} \subset \mathrm{U}(\mathfrak{g})\)。根据 (22)，\(\mathcal{L}\) 是 \(\mathrm{U}(\mathfrak{g})\) 中的一个格。由构造，

\[
\mathcal {U} _ {-}. \mathcal {L} \subset \mathcal {L}\tag{24}
\]

\[
\mathcal {L}. \mathcal {U} _ {+} \subset \mathcal {L}\tag{25}
\]

而引理 3 及 (23) 蕴含

\[
\mathcal {U} _ {0}. \mathcal {L} \subset \mathcal {L}\tag{26}
\]

\[
\mathcal {L}. \mathcal {U} _ {0} \subset \mathcal {L}.\tag{27}
\]

\[
\text {   Let   } \alpha \in \mathrm{B}, n \in \mathbf {N}, r \in \mathbf {N}, \varphi = (\varphi (i)) \in \mathrm{B} ^ {r}, \text {   and   }
\]

\[
(m (1), \dots , m (r)) \in \mathbf {N} ^ {r}.
\]

我们证明

\[
x _ {\alpha} ^ {(n)} y _ {\varphi (1)} ^ {(m (1))} \dots y _ {\varphi (r)} ^ {(m (r))} \in \mathcal {L}\tag{28}
\]

或者，鉴于 (25)，等价地证明

\[
[ x _ {\alpha} ^ {(n)}, y _ {\varphi (1)} ^ {(m (1))} \dots y _ {\varphi (r)} ^ {(m (r))} ] \in \mathcal {L}.\tag{29}
\]

我们对 \(r\) 归纳。要研究的括号是以下各项之和

\[
y _ {\varphi (1)} ^ {(m (1))} \dots y _ {\varphi (k)} ^ {(m (k))} [ x _ {\alpha} ^ {(n)}, y _ {\varphi (k + 1)} ^ {(m (k + 1))} ] y _ {\varphi (k + 2)} ^ {(m (k + 2))} \dots y _ {\varphi (r)} ^ {(m (r))}.\tag{30}
\]

当 \(\alpha \neq \varphi (k + 1)\) 时，\(x_{\alpha}\) 与 \(y_{\varphi (k + 1)}\) 对易，所以 \([x_{\alpha}^{(n)},y_{\varphi (k + 1)}^{(m(k + 1))}] = 0\)。若 \(\alpha = \varphi (k + 1)\)，则由 (17)，表达式 (30) 是以下形式的表达式之和

\[
y _ {\varphi (1)} ^ {(m (1))} \dots y _ {\varphi (k)} ^ {(m (k))} y _ {\varphi (k + 1)} ^ {(m (k + 1) - p)} \binom{q - h}{p} x _ {\alpha} ^ {(n - p)} y _ {\varphi (k + 2)} ^ {(m (k + 2))} \dots y _ {\varphi (r)} ^ {(m (r))}\tag{31}
\]

其中 \(q \in \mathbf{Z}, p \in \mathbf{N} - \{0\}, h \in \mathcal{H}\)。归纳假设连同 (24) 和 (26) 证明表达式 (31) 属于 \(\mathcal{L}\)。于是 (28) 得证。

由 (28)，\(x_{\alpha}^{(n)}\mathcal{U}_{-}\subset \mathcal{L}\)；因而根据 (25) 和 (27)，\(x_{\alpha}^{(n)}\mathcal{L}\subset \mathcal{L}\)，所以 \(\mathcal{U}_{+}.\mathcal{L}\subset \mathcal{L}\)，并且

\[
\mathcal {L}. \mathcal {L} \subset \mathcal {U} _ {-}. \mathcal {U} _ {0}. \mathcal {L} \subset \mathcal {U} _ {-}. \mathcal {L} \subset \mathcal {L}.
\]

因此，\(\mathcal{L}\) 是一个 \(\mathbf{Z}\)-子代数，包含于 \(\mathrm{U}(\mathfrak{g})\)，从而 \(\mathcal{U} = \mathcal{L}\)。若 \(c\) 是 \(\mathrm{U}(\mathfrak{g})\) 的余积，则 \(c(\mathcal{U}) \subset \mathcal{U} \otimes_{\mathbf{Z}} \mathcal{U}\)（第 2 号，命题 1）。令 \(\gamma\) 为 \(\mathrm{U}(\mathfrak{g})\) 的余单位。由于 \(\gamma(x_{\alpha}^{(n)}) = \gamma(y_{\alpha}^{(n)}) = \gamma\left(\binom{h}{n}\right) = 0\) 对 \(n > 0\) 成立，有 \(\gamma(\mathcal{U}) \subset \mathbf{Z}\)。这证明 (i)。另一方面，

\[
\mathcal {U} \cap \mathfrak {h} = \mathcal {L} \cap \mathfrak {h} = \mathcal {U} _ {0} \cap \mathfrak {h} = \mathcal {H}
\]

根据第 4 号的命题 2；类似地，

\[
\mathcal {U} \cap \mathfrak {g} ^ {\alpha} = \mathcal {U} _ {+} \cap \mathfrak {g} ^ {\alpha} = \mathbf {Z} x _ {\alpha}
\]

根据引理 5。这证明 (ii)。

备注 4. 根据第 4 节第 4 号的命题 5，存在唯一的自同构 \(\theta\)，定义在 \(\mathfrak{g}\) 上，并满足 \(\theta(x_{\alpha}) = y_{\alpha}\) 和 \(\theta(y_{\alpha}) = x_{\alpha}\) 对所有 \(\alpha \in \mathrm{B}\) 成立，且 \(\theta(h) = -h\) 对所有 \(h \in \mathfrak{h}\) 成立；有 \(\theta^2 = 1\)。由 \(\mathcal{U}\) 的构造可见，在 \(\mathrm{U}(\mathfrak{g})\) 上延拓 \(\theta\) 所得的自同构保持 \(\mathcal{U}\) 稳定。

推论 1. 置 \(\mathcal{G} = \mathcal{U} \cap \mathfrak{g}\)。则 \(\mathcal{G}\) 是李代数 \(\mathfrak{g}\) 中的一个序，并且在 \(\theta\) 下稳定。有 \(\mathcal{G} = \mathcal{H} + \sum_{\alpha \in \mathrm{R}} (\mathcal{G} \cap \mathfrak{g}^{\alpha})\)。对于所有 \(\alpha \in \mathrm{B}\) 及所有 \(n \in \mathbf{N}\)，映射 \((\operatorname{ad} x_{\alpha})^{n} / n!\)、\((\operatorname{ad} y_{\alpha})^{n} / n!\) 都保持 \(\mathcal{U}\) 和 \(\mathcal{G}\) 稳定。

第一个断言显然。第二个断言由考察类型为 \(Q(R)\) 的分次（定义在 \(U(\mathfrak{g})\) 和 U 上）得到。第三个断言由第 5 号的引理 2 得到。

推论 2. 令 \(w \in \mathrm{W}\)。存在初等自同构 \(\varphi\)，它作用于 \(\mathfrak{g}\)，与 \(\theta\) 对易，保持 \(\mathcal{G}\) 和 \(\mathcal{U}\) 稳定，并延拓 \(w\)。

只需处理 \(w\) 形如 \(s_{\alpha}\)（\(\alpha \in \mathrm{B}\)）的情形。首先注意，\(\operatorname{ad} x_{\alpha}\) 和 \(\operatorname{ad} y_{\alpha}\) 在 \(\mathrm{U}(\mathfrak{g})\) 上局部幂零；换言之，对每个 \(u \in \mathrm{U}(\mathfrak{g})\)，存在一个整数 \(n\)，使得 \((\operatorname{ad} x_{\alpha})^n u = (\operatorname{ad} y_{\alpha})^n u = 0\)。因此可以定义自同构 \(e^{\operatorname{ad} x_{\alpha}} = \sum_{n=0}^{\infty} \frac{1}{n!} (\operatorname{ad} x_{\alpha})^n\) 和 \(e^{\mathrm{ad}y_{\alpha}}\) 作用于 \(\mathrm{U}(\mathfrak{g})\)；立即验证，这些 \(\mathrm{U}(\mathfrak{g})\) 的自同构保持 \(\mathcal{U}\) 稳定。置 \(\varphi_1 = e^{\mathrm{ad}x_\alpha}e^{\mathrm{ad}y_\alpha}e^{\mathrm{ad}x_\alpha}\)、\(\varphi_2 = e^{\mathrm{ad}y_\alpha}e^{\mathrm{ad}x_\alpha}e^{\mathrm{ad}y_\alpha}\)。有 \(\varphi_1|\mathfrak{g} = \varphi_2|\mathfrak{g}\)（第 2 节，第 2 号，公式 (1)），故 \(\varphi_1 = \varphi_2\)。置 \(\varphi_1 = \varphi_2 = \varphi\)。有 \(\theta \varphi \theta^{-1} = \varphi\)，所以 \(\theta\) 与 \(\varphi\) 对易。另一方面，根据第 2 节第 2 号引理 1，有 \(\varphi |\mathfrak{h} = w\)。

推论 3. 令 \(\alpha \in \mathrm{R}\)。若 \(x \in \mathcal{G} \cap \mathfrak{g}^{\alpha}\) 且 \(n \in \mathbf{N}\)，则 \(x^{(n)} \in \mathcal{U}\)，并且 \((\operatorname{ad} x)^n / n!\) 保持 \(\mathcal{G}\) 和 \(\mathcal{U}\) 稳定。

当 \(\alpha \in \mathrm{B}\) 时，由 \(\mathcal{U}\) 的构造及推论 1，这显然成立。在一般情形，存在 \(w \in \mathrm{W}\) 使 \(w(\alpha) \in \mathrm{B}\)（第 VI 章，第 1 节，第 5 号，命题 15）。根据推论 2，存在自同构 \(\varphi\)，作用于 \(\mathfrak{g}\)，保持 \(\mathcal{G}\) 和 \(\mathcal{U}\) 稳定，并将 \(\mathfrak{g}^{\alpha}\) 送到 \(\mathfrak{g}^{w(\alpha)}\)，故由结构迁移得到该推论。

推论 4. 存在一个谢瓦莱系 \((X_{\alpha})_{\alpha \in \mathrm{R}}\)，它位于 \((\mathfrak{g},\mathfrak{h})\) 中（第 2 节，第 4 号，定义 3），并满足 \(X_{\alpha} = x_{\alpha}\) 且 \(X_{-\alpha} = y_{\alpha}\) 对 \(\alpha \in \mathrm{B}\) 成立。对于具有这些性质的任意谢瓦莱系 \((X_{\alpha}^{\prime})_{\alpha \in \mathrm{R}}\)，以及所有 \(\alpha \in \mathrm{R}\)，\(X_{\alpha}^{\prime}\) 都是 \(\mathcal{G} \cap \mathfrak{g}^{\alpha}\) 的一组基。

当 \(\alpha \in \mathrm{B}\) 时，置 \(X_{\alpha} = x_{\alpha}, X_{-\alpha} = y_{\alpha}\)。当 \(\alpha \in \mathrm{R}_{+} - \mathrm{B}\) 时，选取 \(w \in \mathrm{W}\) 使 \(w(\alpha) \in \mathrm{B}\)，并选取自同构 \(\varphi\)，它作用于 \(\mathfrak{g}\)，使得 \(\theta \varphi = \varphi \theta\)、\(\varphi(\mathcal{G}) = \mathcal{G}\)，且 \(\varphi(h) = w^{-1}(h)\) 对 \(h \in \mathfrak{h}\) 成立（推论 2）；置 \(X_{\alpha} = \varphi(x_{w(\alpha)})\)、\(X_{-\alpha} = \varphi(y_{w(\alpha)})\)。则

\[
[ X _ {- \alpha}, X _ {\alpha} ] = \varphi ([ y _ {w (\alpha)}, x _ {w (\alpha)} ]) = \varphi (H _ {w (\alpha)}) = w ^ {- 1} (H _ {w (\alpha)}) = H _ {\alpha}
\]

\[
\theta (X _ {\alpha}) = \theta \varphi (x _ {w (\alpha)}) = \varphi \theta (x _ {w (\alpha)}) = \varphi (y _ {w (\alpha)}) = X _ {- \alpha}
\]

所以 \((X_{\alpha})_{\alpha \in \mathrm{R}}\) 是一个谢瓦莱系。此外，

\[
\mathcal {G} \cap \mathfrak {g} ^ {\alpha} = \varphi (\mathcal {G} \cap \mathfrak {g} ^ {w (\alpha)}) = \varphi (\mathbf {Z} x _ {w (\alpha)}) = \mathbf {Z} X _ {\alpha}\tag{32}
\]

\[
\mathcal {G} \cap \mathfrak {g} ^ {- \alpha} = \varphi (\mathcal {G} \cap \mathfrak {g} ^ {- w (\alpha)}) = \varphi (\mathbf {Z} y _ {w (\alpha)}) = \mathbf {Z} X _ {- \alpha}.\tag{33}
\]

令 \((X'_{\alpha})_{\alpha \in \mathrm{R}}\) 为一个谢瓦莱系，满足 \(X'_{\alpha} = x_{\alpha}, X'_{-\alpha} = y_{\alpha}\) 对 \(\alpha \in \mathrm{B}\) 成立。令 \(\mathrm{S}\) 为满足以下条件的 \(\alpha \in \mathrm{R}\) 所成的集合：\(X'_{\alpha} = \pm X_{\alpha}\)。根据第 2 节第 4 号命题 7，\(\mathrm{S}\) 是根的闭集。由于 \(\mathrm{S} \supset \mathrm{B} \cup (-\mathrm{B})\)，有 \(\mathrm{S} = \mathrm{R}\)（第 VI 章，第 1 节，第 6 号，命题 19）。因此，根据 (32) 和 (33)，有 \(\mathcal{G} \cap \mathfrak{g}^{\alpha} = \mathbf{Z}X'_{\alpha}\) 对所有 \(\alpha \in \mathrm{R}\) 成立。

备注。5) 令 \((X_{\alpha})_{\alpha \in \mathrm{R}}\) 为上面构造的谢瓦莱系。若 \(\alpha, \beta, \alpha + \beta \in \mathrm{R}\)，并置 \([X_{\alpha}, X_{\beta}] = \mathrm{N}_{\alpha, \beta} X_{\alpha + \beta}\)，则有 \([X_{\alpha}, X_{\beta}] \in \mathcal{G} \cap \mathfrak{g}^{\alpha + \beta}\)，从而重新得到 \(\mathrm{N}_{\alpha, \beta} \in \mathbf{Z}\) 这一事实（参见第 2 节，第 4 号，命题 7）。

\begin{enumerate}
\def\labelenumi{\arabic{enumi})}
\setcounter{enumi}{5}
\tightlist
\item
  我们在此过程中还得到谢瓦莱系存在性的新证明（参见第 4 节，第 4 号，命题 5 的推论），且该证明独立于第 2 节的引理 4。
\end{enumerate}

\subsection{CHEVALLEY 序}

设 \((\mathfrak{g},\mathfrak{h})\) 是 \(\mathbf{Q}\) 上的分裂约化李代数，\(\mathrm{R}\) 是其根系。选取：

\begin{enumerate}
\def\labelenumi{\alph{enumi})}
\item
  一个许可格 \(\mathcal{H}\) 位于 \(\mathfrak{h}\) 中（第 6 号，定义 1）；
\item
  对所有 \(\alpha \in \mathrm{R}\)，取一个格 \(\mathcal{G}^{\alpha}\)，位于 \(\mathfrak{g}^{\alpha}\) 中。
\end{enumerate}

置 \(\mathcal{G} = \mathcal{H} \oplus \sum_{\alpha \in \mathrm{R}} \mathcal{G}^{\alpha}\)。这是 \(\mathfrak{g}\) 中的一个格。记 \(\mathcal{U}\) 为一个 \(\mathbf{Z}\)-子代数，包含于 \(\mathrm{U}(\mathfrak{g})\)，它由 \(\binom{h}{n}\)（\(h \in \mathcal{H}, n \in \mathbf{N}\)）和 \(x^{(n)}\)（\(x \in \mathcal{G}^{\alpha}, \alpha \in \mathrm{R}, n \in \mathbf{N}\)）生成。最后，对 \(\alpha \in \mathrm{R}\) 和 \(x \in \mathfrak{g}^{\alpha} - \{0\}\)，置

\[
w _ {\alpha} (x) = (\exp \operatorname{ad} x) (\exp \operatorname{ad} y) (\exp \operatorname{ad} x),
\]

其中 \(y\) 是 \(\mathfrak{g}^{-\alpha}\) 中满足 \([y,x] = H_{\alpha}\) 的唯一元素。采用这些记号：

定理 2. 以下条件等价：

\begin{enumerate}
\def\labelenumi{(\roman{enumi})}
\item
  存在谢瓦莱系 \((X_{\alpha})_{\alpha \in \mathrm{R}}\)，属于 \((\mathfrak{g},\mathfrak{h})\)，并满足 \(\mathcal{G}_{\alpha} = \mathbf{Z}X_{\alpha}\) 对所有 \(\alpha \in \mathrm{R}\) 成立。
\item
  \(\mathcal{U} \cap \mathfrak{g} = \mathcal{G}\)，并且 \([\mathcal{G}^{\alpha}, \mathcal{G}^{-\alpha}] = \mathbf{Z}H_{\alpha}\) 对所有 \(\alpha \in \mathrm{R}\) 成立。
\item
  对所有 \(\alpha \in \mathrm{R}, x \in \mathcal{G}^{\alpha}, n \in \mathbf{N}\)，自同态 \((\operatorname{ad} x)^n / n!\) 作用于 \(\mathfrak{g}\)，将 \(\mathcal{G}\) 映到 \(\mathcal{G}\)，并且 \([\mathcal{G}^{\alpha}, \mathcal{G}^{-\alpha}] = \mathbf{Z}H_{\alpha}\)。
\item
  对所有 \(\alpha \in \mathrm{R}\) 及每个基向量 \(x\)（属于 \(\mathcal{G}^{\alpha}\)），\(w_{\alpha}(x)\) 将 \(\mathcal{G}\) 映到 \(\mathcal{G}\)（即将 \(\mathcal{G}^{\beta}\) 映到 \(\mathcal{G}^{s_{\alpha}(\beta)}\)，其中 \(\beta \in \mathrm{R}\)）。
\item
  \(\Longrightarrow\) (ii)：令 \((X_{\alpha})_{\alpha \in \mathrm{R}}\) 为 \((\mathfrak{g},\mathfrak{h})\) 中满足 \(\mathcal{G}^{\alpha} = \mathbf{Z}X_{\alpha}\) 对所有 \(\alpha \in \mathrm{R}\) 成立的谢瓦莱系，并令 B 为 R 的一组基。对 \(\alpha \in \mathrm{B}\)，置 \(x_{\alpha} = X_{\alpha},y_{\alpha} = X_{-\alpha}\)。令 \(\mathcal{U}'\) 为第 6 号命题 3 根据 \(\mathcal{H}\)、\(x_{\alpha}\) 和 \(y_{\alpha}\) 关联的双序。显然 \(\mathcal{U}' \subset \mathcal{U}\)。根据命题 3 的推论 3 和 4，都有 \(x^{(n)} \in \mathcal{U}'\)，对所有 \(\alpha \in \mathrm{R}\)、\(x \in \mathcal{G}^{\alpha}\) 及 \(n \in \mathbf{N}\) 成立。因此 \(\mathcal{U} = \mathcal{U}'\)，这证明 (ii)。
\item
  \(\Longrightarrow\) (iii)：这由第 5 号的引理 2 立即得到。
\item
  \(\Longrightarrow\) (iv)：令 \(\alpha \in \mathrm{R}\)，并令 \(x\) 为 \(\mathcal{G}^{\alpha}\) 的一个基向量。由于 \([\mathcal{G}^{\alpha},\mathcal{G}^{-\alpha}] = \mathbf{Z}H_{\alpha}\)，唯一的 \(y \in \mathfrak{g}^{-\alpha}\) 满足 \([y,x] = H_{\alpha}\)，并属于 \(\mathcal{G}^{-\alpha}\)。根据 (iii)，\(\exp \mathrm{ad} x\) 和 \(\exp \mathrm{ad} y\) 保持 \(\mathcal{G}\) 稳定，因而 \(w_{\alpha}(x)\) 也如此。
\item
  \(\Longrightarrow\) (i)：令 B 为 R 的一组基。选取一组基 \(x_{\alpha}\)，属于 \(\mathcal{G}^{\alpha}\)，对所有 \(\alpha \in \mathrm{B}\) 成立。令 \(y_{\alpha} \in \mathcal{G}^{-\alpha}\) 满足 \([y_{\alpha}, x_{\alpha}] = H_{\alpha}\)。根据第 1 节第 5 号的公式 (5)，有 \(y_{\alpha} = w_{\alpha}(x_{\alpha}). x_{\alpha}\)，所以根据 (iv)，\(y_{\alpha}\) 是 \(\mathcal{G}^{-\alpha}\) 的一组基。令 \(\mathcal{G}\) 为 \(\mathfrak{g}\) 中由 \(\mathcal{H}\)、\(x_{\alpha}\) 和 \(y_{\alpha}\) 定义的序（第 6 号，命题 3 的推论 1）。于是 \(\mathcal{G}\) 在 \((\operatorname{ad} x_{\alpha})^{n} / n!\)、\((\operatorname{ad} y_{\alpha})^{n} / n!\) 下稳定（同上），因而也在 \(w_{\alpha}(x_{\alpha})\) 下稳定。
\end{enumerate}

现在令 \(\beta \in \mathrm{R}\)。存在 \(\alpha_0, \alpha_1, \ldots, \alpha_r \in \mathrm{B}\)，使得

\[
\beta = s _ {\alpha_ {r}} s _ {\alpha_ {r - 1}} \dots s _ {\alpha_ {1}} (\alpha_ {0})
\]

（第 VI 章，第 1 节，第 5 号，命题 15）。于是，由 (iv)，\(w_{\alpha_r}(x_{\alpha_r}).w_{\alpha_{r-1}}(x_{\alpha_{r-1}}) \ldots w_{\alpha_1}(x_{\alpha_1})\) 将 \(\mathcal{G}^{\alpha_0}\) 映到 \(\mathcal{G}^\beta\)，并且根据前面的结果，将 \(\mathcal{G} \cap \mathfrak{g}^{\alpha_0}\) 映到 \(\mathcal{G} \cap \mathfrak{g}^\beta\)。由于 \(\mathcal{G} \cap \mathfrak{g}^{\alpha_0} = \mathcal{G}^{\alpha_0}\)（命题 3 (ii)），有 \(\mathcal{G} \cap \mathfrak{g}^\beta = \mathcal{G}^\beta\)。因此

\[
\mathcal {G} = \mathcal {H} \oplus \sum_ {\beta \in \mathrm{R}} (\mathcal {G} \cap \mathfrak {g} ^ {\beta}) = \mathcal {H} \oplus \sum_ {\beta \in \mathrm{R}} \mathcal {G} ^ {\beta} = \mathcal {G}
\]

且命题 3 的推论 4 完成证明。

定义 2. 当定理 2 的条件 (i) 至 (iv) 满足时，称 \(\mathcal{G}\) 为 \((\mathfrak{g}, \mathfrak{h})\) 中的谢瓦莱序。

备注。 \((\mathfrak{g},\mathfrak{h})\) 中的谢瓦莱序总是存在。事实上，谢瓦莱序是形如 \(\mathcal{H}\oplus\sum_{\alpha\in \mathrm{R}}\mathbf{Z}X_{\alpha}\) 的集合，其中 \((X_{\alpha})_{\alpha\in \mathrm{R}}\) 是 \((\mathfrak{g},\mathfrak{h})\) 中的一个谢瓦莱系，而 \(\mathcal{H}\) 是 \(\mathfrak{h}\) 中满足

\[
\mathrm{Q} (\mathrm{R} ^ {\vee}) \subset \mathcal {H} \subset \mathrm{P} (\mathrm{R} ^ {\vee}) \oplus \mathfrak {c}
\]

（其中 c 是 g 的中心）。

定理 3. 保持第 7 号开头的记号，并假设 \(\mathcal{G}\) 是 \((\mathfrak{g},\mathfrak{h})\) 中的谢瓦莱序。

\begin{enumerate}
\def\labelenumi{(\roman{enumi})}
\item
  \(\mathcal{U}\) 是 \(\mathrm{U}(\mathfrak{g})\) 中的 b-序。
\item
  令 B 为 \(\mathrm{R}\) 的一组基，并令 \((X_{\alpha})_{\alpha\in\mathrm{B}\cup(-\mathrm{B})}\) 为 \(\mathfrak{g}\) 中满足 \(\mathcal{G}^{\alpha}=\mathbf{Z}X_{\alpha}\) 对 \(\alpha\in\mathrm{B}\cup(-\mathrm{B})\) 成立的一族元素。\(\mathbf{Z}\)-代数 \(\mathcal{U}\) 由 \(\binom{h}{n}\) 和 \(X_{\alpha}^{(n)}\)（\(h\in\mathcal{H},\alpha\in\mathrm{B}\cup(-\mathrm{B}),n\in\mathbf{N}\)）生成。若 \(\mathfrak{g}\) 是半单的且 \(\mathcal{H}=\mathrm{Q}(\mathrm{R}^{\vee})\)，则 \(\mathbf{Z}\)-代数 \(\mathcal{U}\) 由 \(X_{\alpha}^{(n)}\)（\(\alpha\in\mathrm{B}\cup(-\mathrm{B}),n\in\mathbf{N}\)）生成。
\item
  令 B 为 R 的一组基，\(R_{+}\) 为相应的正根集合，\(R_{-} = -R_{+}\)，\(n_{+} = \sum_{\alpha \in R_{+}} g^{\alpha}\)，\(n_{-} = \sum_{\alpha \in R_{-}} g^{\alpha}\)。则
\end{enumerate}

\[
\mathcal {U} = (\mathcal {U} \cap \mathrm{U} (\mathfrak {n} _ {-})). (\mathcal {U} \cap \mathrm{U} (\mathfrak {h})). (\mathcal {U} \cap \mathrm{U} (\mathfrak {n} _ {+})).
\]

令 \((h_i)_{i\in \mathrm{I}}\) 为 \(\mathcal{H}\) 的一组基。对所有 \(\alpha \in \mathrm{R}\)，令 \(X_{\alpha}\) 为 \(\mathcal{G}^{\alpha}\) 的一组基。给集合 \(\mathrm{I}\cup \mathrm{R}\) 赋予一个全序（我们假定 \(\mathrm{I}\cap \mathrm{R} = \emptyset\)）。对 \(\lambda \in \mathrm{I}\cup \mathrm{R}\) 及 \(n\in \mathbf{N}\)，置 \(e_{\lambda}^{\langle n\rangle} = \binom{h_{\lambda}}{n}\)（若 \(\lambda \in \mathrm{I}\)），置 \(e_{\lambda}^{\langle n\rangle} = X_{\lambda}^{(n)}\)（若 \(\lambda \in \mathrm{R}\)）。则乘积 \(\prod_{\lambda \in \mathrm{I}\cup \mathrm{R}}e_{\lambda}^{\langle n_{\lambda}\rangle}\) 当 \((n_{\lambda})\) 属于 \(\mathbf{N}^{\mathrm{I}\cup \mathrm{R}}\) 时，构成 \(\mathbf{Z}\)-模 \(\mathcal{U}\) 的一组基。乘积 \(\prod_{\lambda \in \mathrm{I}}\binom{h_{\lambda}}{n_{\lambda}}\) 当 \((n_{\lambda})\) 属于 \(\mathbf{N}^{\mathrm{I}}\) 时，构成 \(\mathbf{Z}\)-模 \(\mathcal{U}\cap \mathrm{U}(\mathfrak{h})\) 的一组基。乘积 \(\prod_{\lambda \in \mathrm{R}_{+}}X_{\lambda}^{(n_{\lambda})}\) 当 \((n_{\lambda})\) 属于 \(\mathbf{N}^{\mathrm{R}_{+}}\) 时，构成 \(\mathbf{Z}\)-模 \(\mathcal{U}\cap \mathrm{U}(\mathfrak{n}_{+})\) 的一组基。

令 B 和 \((X_{\alpha})_{\alpha\in\mathrm{B}\cup(-\mathrm{B})}\) 如 (ii) 中所述，并满足 \([X_{-\alpha},X_{\alpha}]=H_{\alpha}\)。令 \(\mathcal{U}'\) 为一个 \(\mathbf{Z}\)-子代数，包含于 \(\mathrm{U}(\mathfrak{g})\)，它由 \(\binom{h}{n}\) 和 \(X_{\alpha}^{(n)}\) \((h\in\mathcal{H},\alpha\in\mathrm{B}\cup(-\mathrm{B}),n\in\mathbf{N})\) 生成。在定理 2 的证明中，从 (i) \(\Longrightarrow\) (ii) 的部分，我们已经看到 \(\mathcal{U}'\) 等于 \(\mathcal{U}\)，并且是 \(\mathrm{U}(\mathfrak{g})\) 中的双序。这证明 (i) 及 (ii) 的第一个断言；第二个断言由引理 4 (ii) 得到。断言 (iii) 由定理 1（第 3 号）和命题 3（第 6 号）得到。

\subsection{容许格}

沿用向量空间中采用的术语，若模 M 的自同态 u 存在一组基，使得 u 关于这组基的矩阵为对角矩阵，则称 u 可对角化。

引理 6. 设 M 是有限型自由 \(\mathbf{Z}\)-模，u 是 M 的自同态，v 是自同态 \(u \otimes 1\)，作用于 \(M \otimes_{\mathbf{Z}} \mathbf{Q}\)。假设 \(\binom{v}{n}(M) \subset M\) 对所有 \(n \in \mathbf{N}\) 成立。则 u 可对角化。

\begin{enumerate}
\def\labelenumi{\alph{enumi})}
\item
  对于任意多项式 \(\mathrm{P} \in \mathbf{Q}[\mathrm{T}]\)，若 \(\mathrm{P}(\mathbf{Z}) \subset \mathbf{Z}\)，则有 \(\mathrm{P}(v)(\mathrm{M}) \subset \mathrm{M}\)（第 4 号，命题 2 的推论），所以 det \(\mathrm{P}(v) \in \mathbf{Z}\)。
\item
  记 \(\chi_{v}(t)=t^{d}+\alpha_{1}t^{d-1}+\cdots\) 为 v 的特征多项式。令 \(k\in\mathbf{Z}, n\in\mathbf{N}\)。将 a) 应用于多项式 \(\binom{T-k}{n}\)，可知数
\end{enumerate}

\[
\begin{array}{l} a _ {n} = \det \binom {v - k} {n} = \frac {1}{(n !) ^ {d}} \det (v - k) \det (v - k - 1) \dots \det (v - k - n + 1) \\ = \frac {(- 1) ^ {n}}{(n !) ^ {4}} \chi_ {v} (k) \chi_ {v} (k + 1) \dots \chi_ {v} (k + n - 1) \end{array}
\]

是一个整数。取 \(k - 1 < -\alpha_{1} / d\)。则

\[
\chi_ {v} (k + n - 1) = n ^ {d} + (\alpha_ {1} + (k - 1) d) n ^ {d - 1} + \dots
\]

以及

\[
| a _ {n} | = \frac {| \chi_ {v} (k + n - 1) |}{n ^ {d}} | a _ {n - 1} |;
\]

因此，若 \(a_{n} \neq 0\) 对所有 \(n \in \mathbf{N}\) 成立，则 \(|a_{n}|\) 的序列在 n 充分大时严格递减，这是荒谬的。由此 v 有一个整数特征值 \(\lambda\)。置 \(\mathrm{M}' = \mathrm{Ker}(u - \lambda.1)\)，并置 \(M'' = M/M'\)。则 \(M'\) 是 \(M \otimes_{\mathbf{Z}} \mathbf{Q}\) 的某个向量子空间与 M 的交，因此 \(\mathbf{Z}\)-模 \(M''\) 是有限型无挠模，因而是秩 \textless{} d 的自由模。对 d 归纳，并将归纳假设应用于由 u 在 \(M''\) 上诱导的自同态，可知 v 在 \(\mathbf{Q}\) 的一个代数闭扩张中的所有特征值都是整数。

\begin{enumerate}
\def\labelenumi{\alph{enumi})}
\setcounter{enumi}{2}
\tightlist
\item
  我们证明 \(v\) 可对角化。令 \(\lambda\) 为 \(v\) 的一个特征值，并令 \(x \in \mathrm{M} \otimes_{\mathbf{Z}} \mathbf{Q}\) 满足 \((v - \lambda)^2 x = 0\)。有 \(v(vx - \lambda x) = \lambda(vx - \lambda x)\)，所以
\end{enumerate}

\[
\begin{array}{c} \frac {1}{n !} (v - \lambda - n + 1) (v - \lambda - n + 2) \ldots (v - \lambda - 1) (v - \lambda) x \\ = \frac {(- 1) ^ {n - 1}}{n} (v x - \lambda x). \end{array}
\]

根据 \(a\)，这意味着 \(vx - \lambda x \in n\mathbf{M}\) 对所有 \(n \in \mathbf{N}\) 成立，所以 \((v - \lambda)x = 0\)。

\begin{enumerate}
\def\labelenumi{\alph{enumi})}
\setcounter{enumi}{3}
\tightlist
\item
  令 \(\lambda\) 为 v 的一个特征值，并令 \((\lambda - a, \lambda + b]\) 为 \(\mathbf{Z}\) 中包含 v 的所有特征值的一个区间。考虑多项式
\end{enumerate}

\[
\begin{array}{c} \mathrm{P(T)} = (- 1) ^ {b} \frac {(\mathrm{T} - \lambda - 1) (\mathrm{T} - \lambda - 2) \ldots (\mathrm{T} - \lambda - b)}{b !} \\ \times \frac {(\mathrm{T} - \lambda + 1) (\mathrm{T} - \lambda + 2) \ldots (\mathrm{T} - \lambda + a)}{a !}. \end{array}
\]

有 \(\mathrm{P}(\mathbf{Z})\subset\mathbf{Z},\mathrm{P}(\lambda)=1,\mathrm{P}(\mu)=0\)，其中 \(\mu\in\mathbf{Z}\cap(\lambda-a,\lambda+b)\) 且 \(\mu\neq\lambda\)。根据 a)，\(\mathrm{P}(v)(\mathrm{M})\subset\mathrm{M}\)。根据 c)，\(\mathrm{P}(v)\) 是从 \(M\otimes_{\mathbf{Z}}\mathbf{Q}\) 到对应于 \(\lambda\) 的特征子空间的投影。

证毕。

备注 1. 即使只假设 v 可对角化且特征值为整数，u 也不一定可对角化（例如，取 \(M = \mathbf{Z}^{2}\)，且 \(u(x, y) = (y, x)\) 对所有 \((x, y) \in M\) 成立）。

令 \(\mathfrak{g},\mathfrak{h},\mathrm{R},\mathcal{H},\mathcal{G}^{\alpha},\mathcal{G},\mathcal{U}\) 如第 7 号中所述，并假设 \(\mathcal{G}\) 是 \((\mathfrak{g},\mathfrak{h})\) 中的谢瓦莱序。

定义 3. 令 E 为一个 \(\mathfrak{g}\)-模。若 E 中的格 \(\mathcal{E}\) 满足以下条件，则称其为容许格（相对于 \(\mathcal{G}\)）：

\begin{enumerate}
\def\labelenumi{(\roman{enumi})}
\item
  \(\mathcal{U}\) 将 \(\mathcal{E}\) 映到 \(\mathcal{E}\)；
\item
  \(\mathcal{E}\) 在 \(\binom{h}{n}\) 和 \(x^{(n)}\) 下稳定，对所有 \(\alpha \in \mathrm{R}, x \in \mathcal{G}^{\alpha}, n \in \mathbf{N}, h \in \mathcal{H}\) 成立。
\end{enumerate}

备注。2) 令 \(\rho\) 为 \(\mathfrak{g}\) 在 \(\mathrm{U}(\mathfrak{g})\) 上的伴随表示。令 \(\alpha, x, n, h\) 如上面 (ii) 中所述。由引理 2，有 \(\rho(x^{(n)})\) . \(\mathcal{U} \subset \mathcal{U}\)。另一方面，若 \(p \in \mathbf{N}\)，

\[
\rho \left(\binom{h}{p}\right) x ^ {(n)} = \binom{\operatorname{ad} h}{p} x ^ {(n)} = \binom{n \alpha (h)}{p} x ^ {(n)}
\]

（第 5 号，公式 (13)），所以 \(\rho\left(\binom{h}{p}\right). \mathcal{U} \subset \mathcal{U}\)。这证明 \(\mathcal{U}\) 是 \(\mathrm{U}(\mathfrak{g})\) 中的容许格，从而可知 \(\mathcal{G}\) 是 \(\mathfrak{g}\) 中的容许格（对于伴随表示）。

\begin{enumerate}
\def\labelenumi{\arabic{enumi})}
\setcounter{enumi}{2}
\tightlist
\item
  令 E 为有限维 \(\mathfrak{g}\)-模，\(\mathcal{E}\) 为 E 中的容许格，\(\mathfrak{c}\) 为 \(\mathfrak{g}\) 的中心。根据引理 6，\(\mathfrak{c}\) 的每个元素都在 E 上定义一个可对角化的自同态。因此 E 是半单的（第 I 章，第 6 节，第 5 号，定理 4）。所以，E 是单 \(\mathcal{D}\mathfrak{g}\)-模的直和，且 \(\mathfrak{c}\) 在这些模上诱导数乘变换。根据引理 6，\(\mathcal{E}=\oplus(\mathcal{E}\cap\mathrm{E}^{\lambda})\)，并且对于 E 的所有权 \(\lambda\)，有
\end{enumerate}

\[
\lambda (\mathcal {H}) \subset \mathbf {Z}.
\]

\begin{enumerate}
\def\labelenumi{\arabic{enumi})}
\setcounter{enumi}{3}
\tightlist
\item
  若 \(\mathfrak{g}\) 是半单的且 \(\mathcal{H} = \mathrm{Q}(\mathrm{R}^{\vee})\)，则根据定理 3 (ii)，定义 3 的条件 (i) 和 (ii) 等价于
\end{enumerate}

\begin{enumerate}
\def\labelenumi{(\roman{enumi})}
\setcounter{enumi}{2}
\tightlist
\item
  \(\mathcal{E}\) 在 \(x^{(n)}\) 下稳定，对所有 \(\alpha \in \mathrm{R}, x \in \mathcal{G}^{\alpha}, n \in \mathbf{N}\) 成立。
\end{enumerate}

\begin{enumerate}
\def\labelenumi{\arabic{enumi})}
\setcounter{enumi}{4}
\tightlist
\item
  令 B 为 R 的一组基；在上面的条件 (i) 和 (ii) 中，`` \(\alpha \in R\) '' 可以替换为 `` \(\alpha \in B \cup (-B)\) ''（同上）。
\end{enumerate}

定理 4. 令 \(\mathbf{E}\) 为有限维 \(\mathfrak{g}\)-模。以下条件等价：

\begin{enumerate}
\def\labelenumi{(\roman{enumi})}
\item
  E 有一个容许格；
\item
  \(\mathcal{H}\) 的每个元素都在 \(E\) 上定义一个具有整数特征值的可对角化自同态。
\item
  \(\Longrightarrow\) (ii)：这由备注 3 得到。
\item
  \(\Longrightarrow\) (i)：我们假设条件 (ii) 满足并证明 (i)。根据第 I 章第 6 节第 5 号的定理 4，可以假设 \(\mathfrak{c}\) 的元素在 E 上定义数乘变换，且 E 是单 \(\mathcal{D}\mathfrak{g}\)-模。令 B 为 \(\mathrm{R}\) 的一组基，并令 \(\mathfrak{g} = \mathfrak{n}_{-} \oplus \mathfrak{h} \oplus \mathfrak{n}_{+}\) 为 \(\mathfrak{g}\) 的相应分解。令 \(\lambda\) 为 \(\mathcal{D}\mathfrak{g}\)-模 E 的最高权，并令 \(e \in E^{\lambda} - \{0\}\)。置 \(\mathcal{E} = \mathcal{U}.e\)。显然 \(\mathcal{U}.\mathcal{E} \subset \mathcal{E}\)。由于 E 是单模，\(\mathrm{U}(\mathfrak{g})\) . e = E，从而 \(\mathcal{E}\) 作为 \(\mathbf{Q}\)-向量空间生成 E。对 \(h \in \mathcal{H}\) 和 \(n \in \mathbf{N}\)，有 \(\binom{h}{n} e = \binom{\lambda(h)}{n} e \in \mathbf{Z} e\)，所以
\end{enumerate}

\[
(\mathcal {U} \cap \mathrm{U} (\mathfrak {h})). e = \mathbf {Z}   e.
\]

由于 \(\mathrm{U}(\mathfrak{n}_{+}).e=0\)，根据命题 3 有 \(\mathcal{E}=(\mathcal{U}\cap\mathrm{U}(\mathfrak{n}_{-}))\) .e。现在由定理 3 (iii) 可知，\(\mathcal{E}\) 是有限型 \(\mathbf{Z}\)-模。

推论。若 \(\mathfrak{g}\) 是半单的且 \(\mathcal{H}=\mathrm{Q}(\mathrm{R}^{\vee})\)，则每个有限维 \(\mathfrak{g}\)-模都有一个容许格。

\section{经典可分裂单李代数}

在本段中，我们将对每一种经典可分裂单李代数明确描述：

\begin{enumerate}
\def\labelenumi{(\Roman{enumi})}
\item
  这种类型的一个代数、其维数及其可分裂嘉当子代数；
\item
  其余根；
\item
  其 Borel 子代数及其抛物子代数；
\item
  其基本单表示；
\item
  其基本单表示中为正交型或辛型的那些表示；
\item
  不变多项式函数的代数；
\item
  群 \(\operatorname{Aut} \mathfrak{g}\)、\(\operatorname{Aut}_0 \mathfrak{g}\)、\(\operatorname{Aut}_e \mathfrak{g}\) 的某些性质；
\item
  Killing 型在嘉当子代数上的限制；
\item
  谢瓦莱序。
\end{enumerate}

\subsection{\texorpdfstring{\(\mathbf{A}_l\) 型代数 ( \(l \geq 1\) )}{\textbackslash mathbf\{A\}\_l 型代数 ( l \textbackslash geq 1 )}}

\begin{enumerate}
\def\labelenumi{(\Roman{enumi})}
\tightlist
\item
  令 V 为 k 上维数为 \(l+1\) 的向量空间，并令 \(\mathfrak{g}\) 为 V 的迹为零的自同态代数 \(\mathfrak{sl}(V)\)。令 \((e_{i})_{1\leq i\leq l+1}\) 为 V 的一组基；将 \(\mathfrak{g}\) 的元素关联到其关于这组基的矩阵的映射，是 \(\mathfrak{g}\) 与迹为零矩阵代数 \(\mathfrak{sl}(l+1,k)\) 之间的一个识别。我们知道 \(\mathfrak{g}\) 是半单的（第 I 章，第 6 节，第 7 号，命题 8）。
\end{enumerate}

回忆（代数，第 II 章，第 10 节，第 3 号），\(E_{ij}\) 表示矩阵 \((\alpha_{mp})\)，其中 \(\alpha_{ij} = 1\)，且 \(\alpha_{mp} = 0\) 对 \((m,p) \neq (i,j)\) 成立。矩阵

\[
\begin{array}{l l} E _ {i j} & (1 \leq i, j \leq l + 1, i \neq j) \\ E _ {i, i} - E _ {i + 1, i + 1} & (1 \leq i \leq l) \end{array}
\]

构成 \(\mathfrak{g}\) 的一组基。因此

\[
\dim \mathfrak {g} = l (l + 2).
\]

令 \(\hat{\mathfrak{h}}\) 为 \(\mathfrak{gl}(l + 1,k)\) 的对角元素集合；序列 \((E_{ii})_{1\leq i\leq l + 1}\) 是向量空间 \(\hat{\mathfrak{h}}\) 的一组基；令 \((\hat{\varepsilon}_i)_{1\leq i\leq l + 1}\) 为 \(\hat{\mathfrak{h}}^*\) 的一组基，且与 \((E_{ii})_{1\leq i\leq l + 1}\) 对偶。对所有 \(h\in \hat{\mathfrak{h}}\)，

\[
[ h, E _ {i j} ] = (\hat {\varepsilon} _ {i} (h) - \hat {\varepsilon} _ {j} (h)) E _ {i j}\tag{1}
\]

根据第 I 章，第 1 节，第 2 号的公式 (5)。令 \(\mathfrak{h}\) 为 \(\hat{\mathfrak{h}}\) 中迹为零的元素集合，并置 \(\varepsilon_i = \hat{\varepsilon}_i|\mathfrak{h}\)。则 \(\mathfrak{h}\) 是 \(\mathfrak{g}\) 的嘉当子代数（第 VII 章，第 2 节，第 1 号，例 4）。关系 (1) 证明此嘉当子代数是可分裂的，并且 \((\mathfrak{g},\mathfrak{h})\) 的根为 \(\varepsilon_i - \varepsilon_j\)（\(i\neq j\)）。令 \(\hat{\mathfrak{h}}_0^*\) 为 \(\hat{\mathfrak{h}}^*\) 中关于 \((\hat{\varepsilon}_i)\) 的坐标之和为零的元素集合。映射 \(\lambda \mapsto \lambda |\mathfrak{h}\) 从 \(\hat{\mathfrak{h}}_0^*\) 到 \(\mathfrak{h}^*\) 是双射。因此，\((\mathfrak{g},\mathfrak{h})\) 的根系 R 是 \(A_l\) 型（第 VI 章，第 4 节，第 7 号）。所以 \(\mathfrak{g}\) 是单的（第 3 节，第 2 号，命题 6 的推论 1）。因此，\(\mathfrak{g}\) 是 \(A_l\) 型的可分裂单李代数。

每个可分裂嘉当子代数 \(\mathfrak{h}^{\prime}\) 都是 \(\mathfrak{g}\) 中的 \(\mathfrak{h}\) 在某个初等自同构下的变换（\(\S3\)，第 3 号，命题 10 的推论）。由于 \(\operatorname{Aut}_{e}\mathfrak{g}\) 是由 \(x \mapsto sxs^{-1}\) 给出的 \(\mathfrak{g}\) 的自同构集合，其中 \(s \in \mathbf{SL}(V)\)（\(\S3\)，第 1 号，备注 2；另见 (VII)），存在 V 的一组基 \(\beta\)，使得 \(\mathfrak{h}^{\prime}\) 等于 \(\mathfrak{h}_{\beta}\)，即由 \(\mathfrak{g}\) 中关于基 \(\beta\) 的矩阵为对角矩阵的元素组成的集合。由于 \(\mathfrak{h}_{\beta}\) 含有一个具有互异特征值的元素，V 中在 \(\mathfrak{h}_{\beta}\) 的元素作用下稳定的向量子空间，只有由 \(\beta\) 的子集生成的那些。由此，映射 \(\beta \mapsto \mathfrak{h}_{\beta}\) 通过取商，诱导出从 V 分解为 \(l + 1\) 个一维子空间的分解集合到 \(\mathfrak{g}\) 的可分裂嘉当子代数集合的双射。

\begin{enumerate}
\def\labelenumi{(\Roman{enumi})}
\setcounter{enumi}{1}
\tightlist
\item
  令 \(\alpha = \varepsilon_{i} - \varepsilon_{j}\)（\(i \neq j\)）为一个根。有 \(\mathfrak{g}^{\alpha} = kE_{ij}\)。由于
\end{enumerate}

\[
[ E _ {i j}, E _ {j i} ] = E _ {i i} - E _ {j j}
\]

并且由于 \(\alpha(E_{ii} - E_{jj}) = 2\)，有（第 2 节，第 2 号，定理 1 (ii)）

\[
H _ {\alpha} = E _ {i i} - E _ {j j}.
\]

\begin{enumerate}
\def\labelenumi{(\Roman{enumi})}
\setcounter{enumi}{2}
\tightlist
\item
  置 \(\alpha_{1} = \varepsilon_{1} - \varepsilon_{2},\alpha_{2} = \varepsilon_{2} - \varepsilon_{3},\ldots ,\alpha_{l} = \varepsilon_{l} - \varepsilon_{l + 1}\)。根据第 VI 章，第 4 节，第 7 号.I，\((\alpha_{1},\dots,\alpha_{l})\) 是 R 的一组基 B；相对于 B 的正根是 \(\varepsilon_i - \varepsilon_j\)，其中 \(i < j\)。相应的 Borel 子代数 \(\mathfrak{b}\) 是迹为零的上三角矩阵集合。
\end{enumerate}

V 中的一个旗是 V 的向量子空间的一个集合，它不同于 \(\{0\}\) 和 V，并按包含关系全序。按包含关系给 V 的旗集合赋序。极大旗是集合 \(\{W_{1},\ldots,W_{l}\}\)，其中 \(W_{i}\) 是 i 维向量子空间，且

\[
\mathrm{W} _ {1} \subset \dots \subset \mathrm{W} _ {l}.
\]

例如，若 \(V_{i}\) 表示由 \(e_{1},\ldots,e_{i}\) 生成的 V 的子空间，则 \(\{V_{1},\ldots,V_{l}\}\) 是一个极大旗。

立即可见，\(\mathfrak{b}\) 是 \(\mathfrak{g}\) 中保持极大旗 \(\{V_{1},\ldots,V_{l}\}\) 的各元素不变的元素集合。反之，由于 \(\mathfrak{b}\) 包含 \(\mathfrak{h}\) 以及矩阵 \(E_{ij}\)，其中 i\textless j，可见 \(V_{i}\) 是在 \(\mathfrak{b}\) 下稳定的非平凡向量子空间中仅有的那些。

现在令 \(\delta\) 为 V 中的一个极大旗。由前述可知，\(\mathfrak{b}_{\delta}\)，即 \(\mathfrak{g}\) 中保持 \(\delta\) 的所有元素不变的元素集合，是 \(\mathfrak{g}\) 的一个 Borel 子代数。由于 \(\mathfrak{g}\) 的每个 Borel 子代数都是 \(\mathfrak{b}\) 在某个初等自同构下的变换，可知映射 \(\delta \mapsto \mathfrak{b}_{\delta}\) 是从极大旗集合到 \(\mathfrak{g}\) 的 Borel 子代数集合的双射。

令 \(\beta\) 为 V 的一组基。根据 (I) 及前述，包含 \(\mathfrak{h}_{\beta}\) 的 Borel 子代数正是与极大旗相对应的那些子代数，其中每个元素都由 \(\beta\) 的某个子集生成。这些旗与 \(\beta\) 上的全序一一对应，具体如下：该全序 \(\omega\) 对应于 \(\beta\) 上的旗 \(\{\mathrm{W}_1,\dots ,\mathrm{W}_l\}\)，其中 \(\mathrm{W}_i\) 是由前 \(i\) 个属于 \(\beta\) 的元素生成的向量子空间，对应的全序为 \(\omega\)。由于 \((l + 1)!\) 个全序定义在 \(\beta\) 上，我们重新得到 \((l + 1)!\) 个 \((\mathfrak{sl}(\mathrm{V}),\mathfrak{h}_{\beta})\) 中的 Borel 子代数这一事实（\(\S 3\)，第 3 号，备注）。

令 \(\gamma\) 为 V 中的一个旗。由于 \(\gamma\) 包含于某个极大旗，\(\mathfrak{p}_{\gamma}\)，即 \(\mathfrak{g}\) 中保持 \(\gamma\) 的各元素不变的元素集合，是 \(\mathfrak{g}\) 的一个抛物子代数。我们证明，在 \(\mathfrak{p}_{\gamma}\) 下稳定的非平凡向量子空间只有 \(\gamma\) 的元素。为此可设 \(\gamma = \{V_{i_1}, \ldots, V_{i_q}\}\)，其中 \(1 \leq i_1 < \cdots < i_q \leq l\)。置 \(i_0 = 0, i_{q+1} = l + 1\)。非空区间

\[
\left[ i _ {0} + 1, i _ {1} \right], \left[ i _ {1} + 1, i _ {2} \right], \dots , \left[ i _ {q} + 1, i _ {q + 1} \right]
\]

构成 \(\{1, \ldots, l + 1\}\) 的一个划分，从而任意 \(l + 1\) 阶方阵都可写成分块矩阵 \((X_{ab})_{1 \leq a, b \leq q + 1}\)。代数 \(\mathfrak{p}_{\gamma}\) 就是集合 \(\mathfrak{p}_{i_1,\dots,i_q}\)，它由元素 \((X_{ab})_{1\leq a,b\leq q + 1}\) 中属于 \(\mathfrak{sl}(l + 1,k)\) 且满足 \(X_{ab} = 0\)（当 \(a > b\)）的那些组成。由于 \(\mathfrak{p}_{i_1,\dots,i_q}\supset \mathfrak{b}\)，在 \(\mathfrak{p}_{i_1,\dots,i_q}\) 下稳定的非平凡向量子空间必为某个 \(\mathrm{V}_i\)；若 \(i_k < i < i_{k + 1}\)，代数 \(\mathfrak{p}_{i_1,\dots,i_q}\) 包含 \(E_{i_{k + 1},i}\)，而 \(\mathrm{V}_i\) 不稳定，故断言成立。

因此，\(2^{l}\) 个包含在极大旗 \(\{V_{1},\ldots,V_{l}\}\) 中的旗产生 \(2^{l}\) 个不同的抛物子代数，且这些子代数包含 \(\mathfrak{b}\)；由于共有 \(2^{l}\) 个包含 \(\mathfrak{b}\) 的抛物子代数（\(\S3\)，第 4 号，备注），可知映射 \(\gamma\mapsto \mathfrak{p}_{\gamma}\) 是从 V 的旗集合到 \(\mathfrak{g}\) 的抛物子代数集合的双射。此外，\(\mathfrak{p}_{\gamma}\supset \mathfrak{p}_{\gamma'}\) 当且仅当 \(\gamma\subset\gamma'\)。

回忆抛物子代数 \(\mathfrak{p} = \mathfrak{p}_{i_1, \ldots, i_q}\)（\(1 \leq i_1 < \cdots < i_q \leq l\)）。令 \(\mathfrak{s}\)（分别为 \(\mathfrak{n}\)）为由矩阵 \((X_{ab})_{1 \leq a, b \leq q+1}\) 构成的集合，这些矩阵属于 \(\mathfrak{sl}(l+1, k)\)，并满足 \(X_{ab} = 0\)：前者当 \(a \neq b\)，后者当 \(a \geq b\)。根据第 3 节第 4 号的命题 13，有 \(\mathfrak{p} = \mathfrak{s} \oplus \mathfrak{n}\)，子代数 \(\mathfrak{s}\) 在 \(\mathfrak{g}\) 中是约化的，而 \(\mathfrak{n}\) 既是 \(\mathfrak{p}\) 的最大幂零理想，也是其幂零根基。

\begin{enumerate}
\def\labelenumi{(\Roman{enumi})}
\setcounter{enumi}{3}
\tightlist
\item
  对 \(r = 1, 2, \ldots, l\)，令 \(\varpi_{r} = \varepsilon_{1} + \cdots + \varepsilon_{r}\)。有 \(\varpi_{i}(H_{\alpha_{j}}) = \delta_{ij}\)，所以 \(\varpi_{r}\) 是对应于 \(\alpha_{r}\) 的基本权。
\end{enumerate}

令 \(\sigma\) 为 \(\mathfrak{g}\) 在 V 上的恒等表示。\(\bigwedge^{r}\sigma\) 是 \(\sigma\) 在 \(\mathrm{E} = \bigwedge^{r}(\mathrm{V})\) 上的表示。令 \((e_1,\ldots ,e_{l + 1})\) 为 V 的所选基。这些 \(e_{i_1}\wedge \dots \wedge e_{i_r}\)（其中 \(i_1 < \dots < i_r\)）构成 E 的一组基。若 \(h\in \mathfrak{h}\)，

\[
(\bigwedge^ {r} \sigma) (h). e _ {i _ {1}} \wedge \dots \wedge e _ {i _ {r}} = (\varepsilon_ {i _ {1}} + \dots + \varepsilon_ {i _ {r}}) (h) e _ {i _ {1}} \wedge \dots \wedge e _ {i _ {r}}.
\]

因此，每个权的重数都是 1，\(\varpi_{r}\) 是 \(\bigwedge^{r}\sigma\) 的一个权，而其他每个权都形如 \(\varpi_{r}-\mu\)，其中 \(\mu\) 是正根权。所以，\(\varpi_{r}\) 是 \(\bigwedge^{r}\sigma\) 的最高权，而 \(e_{1}\wedge\cdots\wedge e_{r}\) 是一个本原元。根据第 VI 章，第 4 节，第 7 号.IX，Weyl 群可识别为集合

\[
\{\varepsilon_ {1}, \dots , \varepsilon_ {l + 1} \}.
\]

因此，\(\varpi_{r}\) 在 Weyl 群下的轨道包含所有 \(\varepsilon_{i_{1}} + \cdots + \varepsilon_{i_{r}}\)，其中 \(i_{1} < \cdots < i_{r}\)。由本原元 \(e_{1} \wedge \cdots \wedge e_{r}\) 生成的单子模因而以所有 \(\varepsilon_{i_{1}} + \cdots + \varepsilon_{i_{r}}\) 为权，从而等于 E。因此，\(\bigwedge^{r} \sigma\) 是最高权为 \(\varpi_{r}\) 的不可约表示。

因此，表示 \(\bigwedge^{r}\sigma\)（\(1 \leq r \leq l\)）是基本表示。有 \(\dim(\bigwedge^{r}\sigma) = \binom{l+1}{r}\)。

\begin{enumerate}
\def\labelenumi{(\Alph{enumi})}
\setcounter{enumi}{21}
\tightlist
\item
  有 \(w_0(\alpha_1) = -\alpha_l, w_0(\alpha_2) = -\alpha_{l-1}, \ldots\)（第 VI 章，第 4 节，第 7 号，XI），因此
\end{enumerate}

\[
- w _ {0} (\varpi_ {1}) = \varpi_ {l}, - w _ {0} (\varpi_ {2}) = \varpi_ {l - 1}, \ldots .
\]

令

\[
\omega = n _ {1} \varpi_ {1} + \dots + n _ {l} \varpi_ {l} \quad (n _ {1}, \ldots , n _ {l} \in \mathbf {N})
\]

为一个支配权。具有最高权 \(\omega\) 的单表示是正交型或辛型，当且仅当

\[
n _ {1} = n _ {l}, \quad n _ {2} = n _ {l - 1}, \ldots
\]

（第 7 节，第 5 号，命题 12）。特别地，若 \(l\) 为偶数，则 \(\mathfrak{sl}(l + 1,k)\) 的基本表示没有一个是正交型或辛型。若 \(l\) 为奇数，则表示 \(\bigwedge^i\sigma\)（其中 \(i\neq (l + 1) / 2\)）既非正交型也非辛型；根据第 VI 章，第 4 节，第 7 号.VI，\(\varpi_{(l + 1) / 2}\) 关于 \((\alpha_{1},\ldots ,\alpha_{l})\) 的坐标之和为

\[
\begin{array}{r l} \frac {1}{l + 1} \left[ \frac {l + 1}{2} \left(1 + 2 + \dots + \frac {l - 1}{2}\right) + \frac {l + 1}{2} \left(1 + 2 + \dots + \frac {l + 1}{2}\right) \right] & = 1 + 2 + \dots + \frac {l - 1}{2} + \frac {l + 1}{4} \end{array}
\]

所以，\(\bigwedge^{(l + 1) / 2}\sigma\) 在 \(l\equiv -1\)（mod. 4）时是正交型的；在 \(l\equiv 1\)（mod. 4）时是辛型的（第 7 节，第 5 号，命题 12）。这一结果还可如下精确化。选取一个非零元素 \(e\)，属于 \(\bigwedge^{l + 1}(\mathrm{V})\)。外代数 \(\mathrm{V}\) 中的乘法定义了从

\[
\bigwedge^ {(l + 1) / 2} (\mathrm{V}) \times \bigwedge^ {(l + 1) / 2} (\mathrm{V})
\]

到 \(\bigwedge^{l + 1}(\mathrm{V})\) 的双线性映射，该映射可写成 \((u,v)\mapsto \varPhi (u,v)e\)，其中 \(\varPhi\) 是 \(\bigwedge^{(l + 1) / 2}(\mathrm{V})\) 上的双线性型。立即可验证，\(\varPhi\) 非零且在 \(\mathfrak{g}\) 下不变（因而非退化）；当 \((l + 1) / 2\) 为偶数时它对称，当 \((l + 1) / 2\) 为奇数时它交错。

\begin{enumerate}
\def\labelenumi{(\Roman{enumi})}
\setcounter{enumi}{5}
\tightlist
\item
  对所有 \(x \in \mathfrak{g}\)，\(\sigma(x) = x\) 的特征多项式可写成
\end{enumerate}

\[
\mathrm{T} ^ {l + 1} + f _ {2} (x) \mathrm{T} ^ {l - 1} + f _ {3} (x) \mathrm{T} ^ {l - 2} + \dots + f _ {l + 1} (x)
\]

其中 \(f_{2}, \ldots, f_{l+1}\) 是在 \(\mathfrak{g}\) 下不变的多项式函数（第 8 节，第 3 号，引理 2）。

若 \(x = \xi_{1}E_{11} + \cdots + \xi_{l+1}E_{l+1 l+1} \in \mathfrak{h}\)，则 \(f_{i}(x)\) 除了符号外，正是 \(\xi_{1}, \ldots, \xi_{l+1}\) 的次数为 \(2, \ldots, l + 1\) 的初等对称函数。因此，根据第 VI 章，第 4 节，第 7 号.IX，\(f_{i}|_{\mathfrak{h}}\) 生成 \(\mathbf{S}(\mathfrak{h}^{*})\) 中在 Weyl 群下不变的元素所成的代数，并且代数无关。所以（第 8 节，第 3 号，命题 3），\(f_{2}, f_{3}, \ldots, f_{l+1}\) 生成在 \(\mathfrak{g}\) 下不变的多项式函数所成的代数，并且代数无关。

\begin{enumerate}
\def\labelenumi{(\Roman{enumi})}
\setcounter{enumi}{6}
\tightlist
\item
  对所有 \(g \in \mathbf{GL}(l + 1, k)\)，令 \(\varphi_k(g) = \varphi(g)\) 为由 \(x \mapsto gxg^{-1}\) 给出的 \(\mathfrak{g}\) 的自同构。则 \(\varphi\) 是从 \(\mathbf{GL}(l + 1, k)\) 到 \(\operatorname{Aut}(\mathfrak{g})\) 的同态。有
\end{enumerate}

\[
\varphi (\mathbf {S L} (l + 1, k)) = \operatorname{Aut} _ {e} (\mathfrak {g})
\]

（第 VII 章，第 3 节，第 1 号，备注 2）。令 \(\bar{k}\) 为 \(k\) 的代数闭包。有

\[
\mathbf {G L} (l + 1, \bar {k}) = \bar {k} ^ {*}. \mathbf {S L} (l + 1, \bar {k}),
\]

所以 \(\varphi_{\bar{k}}(\mathbf{GL}(l + 1,\bar{k})) = \varphi_{\bar{k}}(\mathbf{SL}(l + 1,\bar{k})) = \mathrm{Aut}_e(\mathfrak{g}\otimes_k\bar{k})\)；由此 \(\varphi (\mathbf{GL}(l + 1,k))\subset \mathrm{Aut}_0(\mathfrak{g})\)。另一方面，\(\mathrm{Aut}_0(\mathfrak{g})\subset \varphi (\mathbf{GL}(l + 1,k))\) 由将第 7 节第 1 号的命题 2 应用于 \(\mathfrak{g}\) 的恒等表示而得。因此，

\[
\operatorname{Aut} _ {0} (\mathfrak {g}) = \varphi (\mathbf {G L} (l + 1, k)).
\]

 \(\varphi\) 的核是 \(\mathbf{GL}(l+1,k)\) 中与每个 \(l+1\) 阶矩阵对易的元素集合，即可逆标量矩阵的集合 \(k^{*}\)。因此，\(\operatorname{Aut}_{0}(\mathfrak{g})\) 可识别为群 \(\mathbf{GL}(l+1,k)/k^{*}=\mathbf{PGL}(l+1,k)\)。\(\varphi'=\varphi|\mathbf{SL}(l+1,k)\) 的核是 \(\mu_{l+1}(k)\)，其中 \(\mu_{l+1}(k)\) 表示 k 中的 \((l+1)\) 次单位根集合。因此，\(\operatorname{Aut}_{e}(\mathfrak{g})\) 可识别为群 \(\mathbf{SL}(l+1,k)/\mu_{l+1}(k)=\mathbf{PSL}(l+1,k)\)。另一方面，有正合序列

\[
1 \longrightarrow \mathbf {S L} (l + 1, k) \longrightarrow \mathbf {G L} (l + 1, k) \stackrel {\det} {\longrightarrow} k ^ {*} \longrightarrow 1
\]

且 \(k^*\) 在 det 下的像为 \(k^{*l + 1}\)。由此得到典则同构

\[
\begin{array}{c} \operatorname{Aut} _ {0} (\mathfrak {g}) / \operatorname{Aut} _ {e} (\mathfrak {g}) \longrightarrow \mathbf {P G L} (l + 1, k) / \mathbf {P S L} (l + 1, k) \\ \qquad \qquad \qquad \qquad \qquad \longrightarrow \mathbf {G L} (l + 1, k) / k ^ {*}. \mathbf {S L} (l + 1, k) \longrightarrow k ^ {*} / k ^ {*   l + 1}. \end{array}
\]

若 \(k = \mathbf{R}\)，可见 \(\mathrm{Aut}_0(\mathfrak{g}) = \mathrm{Aut}_e(\mathfrak{g})\) 当 \(l + 1\) 为奇数时成立，而 \(\mathrm{Aut}_0(\mathfrak{g}) / \mathrm{Aut}_e(\mathfrak{g})\) 同构于 \(\mathbf{Z}/2\mathbf{Z}\) 当 \(l + 1\) 为偶数时。

采用第 5 节的记号，\(f(\mathrm{T}_{\mathbf{Q}})\) 是 \(\mathfrak{g}\) 的在 \(\mathfrak{h}\) 上诱导恒等映射的自同构集合，因而等于 \(\varphi(\mathrm{D})\)，其中 D 是 \(\mathbf{GL}(l+1,k)\) 的对角元素集合（第 5 节，命题 4）。令 \(D'\) 为 \(\mathbf{SL}(l+1,k)\) 的对角元素集合。根据第 5 节的命题 3 以及 \(\operatorname{Aut}_{e}(\mathfrak{g})\) 的确定，有 \(f(q(\mathrm{T}_{\mathrm{P}})) \subset \varphi(\mathrm{D}')\)。我们证明 \(f(q(\mathrm{T}_{\mathrm{P}})) = \varphi(\mathrm{D}')\)。令

\[
d = \left( \begin{array}{c c c} \lambda_ {1} & & 0 \\ & \ddots & \\ 0 & & \lambda_ {l + 1} \end{array} \right)
\]

为 \(D'\) 的一个元素。存在 \(\zeta \in \operatorname{Hom}(\mathbf{Q}(\mathbf{R}), k^{*}) = \mathrm{T}_{\mathbf{Q}}\)，使得对所有 i 和 j 都有 \(\zeta(\varepsilon_{i} - \varepsilon_{j}) = \lambda_{i}\lambda_{j}^{-1}\)。容易验证 \(f(\zeta) = \varphi(d)\)。根据第 VI 章，第 4 节，第 7 号.VIII，\(\mathrm{P}(\mathrm{R})\) 由 \(\mathrm{Q}(\mathrm{R})\) 和元素 \(\varepsilon = \varepsilon_{1}\) 生成，后者在 \(\mathrm{P}(\mathrm{R})/\mathrm{Q}(\mathrm{R})\) 中的像的阶为 \(l + 1\)；但是

\[
\begin{array}{c} \zeta ((l + 1) \varepsilon) = \zeta ((\varepsilon_ {1} - \varepsilon_ {2}) + (\varepsilon_ {1} - \varepsilon_ {3}) + \dots + (\varepsilon_ {1} - \varepsilon_ {l + 1})) \\ = \lambda_ {1} ^ {l} \lambda_ {2} ^ {- 1} \lambda_ {3} ^ {- 1} \ldots \lambda_ {l + 1} ^ {- 1} = \lambda_ {1} ^ {l + 1} \end{array}
\]

所以 \(\zeta\) 延拓为从 \(\mathrm{P}(\mathrm{R})\) 到 \(k^{*}\) 的同态。这证明 \(\zeta \in q(\mathrm{T}_{\mathrm{P}})\)，从而 \(\varphi(d) \in f(q(\mathrm{T}_{\mathrm{P}}))\)。

回忆（第 5 节，第 3 号，命题 5 的推论 2），\(\mathrm{Aut}(\mathfrak{g}) = \mathrm{Aut}_0(\mathfrak{g})\) 当 \(l = 1\) 时，且 \(\mathrm{Aut}(\mathfrak{g}) / \mathrm{Aut}_0(\mathfrak{g})\) 同构于 \(\mathbf{Z}/2\mathbf{Z}\) 当 \(l \geq 2\) 时。映射 \(\theta : x \mapsto -^t x\) 是 \(\mathfrak{sl}(l + 1, k)\) 的一个自同构，且 \(a_0 = \theta|\mathfrak{h} \notin \mathrm{W}\) 当 \(l \geq 2\) 时（第 VI 章，第 4 节，第 7 号.XI），所以 \(a_0\) 在 \(\mathrm{Aut}(\mathfrak{g}) / \mathrm{Aut}_0(\mathfrak{g})\) 中的类是该群的非平凡元素（第 5 节，第 2 号，命题 4）。

\begin{enumerate}
\def\labelenumi{(\Roman{enumi})}
\setcounter{enumi}{7}
\tightlist
\item
  Killing 型在 \(\mathfrak{h}\) 上的限制为
\end{enumerate}

\[
\begin{array}{l} \Phi (\xi_ {1} E _ {1 1} + \dots + \xi_ {l + 1} E _ {l + 1, l + 1}, \xi_ {1} ^ {\prime} E _ {1 1} + \dots + \xi_ {l + 1} ^ {\prime} E _ {l + 1, l + 1}) \\ = \sum_ {i \neq j} (\xi_ {i} - \xi_ {j}) (\xi_ {i} ^ {\prime} - \xi_ {j} ^ {\prime}) = \sum_ {i, j} (\xi_ {i} - \xi_ {j}) (\xi_ {i} ^ {\prime} - \xi_ {j} ^ {\prime}) \\ = (l + 1) \sum_ {i} \xi_ {i} \xi_ {i} ^ {\prime} + (l + 1) \sum_ {j} \xi_ {j} \xi_ {j} ^ {\prime} - 2 \left(\sum_ {i} \xi_ {i}\right) \left(\sum_ {j} \xi_ {j} ^ {\prime}\right) \\ = 2 (l + 1) \sum_ {i} \xi_ {i} \xi_ {i} ^ {\prime}. \end{array}
\]

\begin{enumerate}
\def\labelenumi{(\Roman{enumi})}
\setcounter{enumi}{8}
\tightlist
\item
  对 \(1 \leq i < j \leq l + 1\)，置
\end{enumerate}

\[
X _ {\varepsilon_ {i} - \varepsilon_ {j}} = E _ {i j} X _ {\varepsilon_ {j} - \varepsilon_ {i}} = - E _ {j i}.
\]

那么，对所有 \(\alpha \in \mathrm{R}\)，有 \([X_{\alpha}, X_{-\alpha}] = -H_{\alpha}\) 和 \(\theta(X_{\alpha}) = X_{-\alpha}\)（其中 \(\theta\) 是在（VII）中引入的自同构 \(x \mapsto -^{t} x\)）。因此，\((X_{\alpha})_{\alpha \in \mathrm{R}}\) 是 \((\mathfrak{g}, \mathfrak{h})\) 中的一个谢瓦莱系。

取 \(k = \mathbf{Q}\)。\(\mathfrak{h}\) 中的许可格（第 12 节，第 6 号，定义 1）正是介于由 \(\mathbf{Z}\)-模 \(\mathrm{Q}(\mathrm{R}^{\vee})\) 生成的 \(E_{ii} - E_{i+1,i+1}\)（即由属于 \(\mathfrak{sl}(l+1,\mathbf{Z})\) 的对角矩阵组成）与由 \(\mathbf{Z}\)-模 \(\mathrm{P}(\mathrm{R}^{\vee})\) 生成的 \(\mathrm{Q}(\mathrm{R}^{\vee})\) 及 \(E_{11} - (l+1)^{-1}\sum E_{ii}\)（第 VI 章，第 4 节，第 7 号.VIII，即由形如 \(x + (l+1)^{-1}a.1\) 的迹为零对角矩阵组成，其中 \(x\) 的元素为整数且 \(a\in\mathbf{Z}\)）之间的那些格。因此，\(\mathfrak{sl}(l+1,\mathbf{Z})\) 是 \((\mathfrak{g},\mathfrak{h})\) 中与许可格 \(\mathrm{Q}(\mathrm{R}^{\vee})\) 及谢瓦莱系 \((X_{\alpha})\) 相联系的谢瓦莱序。容易验证，\(\bigwedge^{r}\mathbf{Z}^{l+1}\) 是 \(\bigwedge^{r}\mathbf{Q}^{l+1}\) 中相对于 \(\mathfrak{sl}(l+1,\mathbf{Z})\) 的一个容许格（第 12 节，第 8 号，定义 3）。

另一方面，\(\mathfrak{gl}(l+1,\mathbf{Z})\) 是分裂约化代数 \(\mathfrak{gl}(l+1,\mathbf{Q})\) 中的谢瓦莱序；它在 \(\mathfrak{sl}(l+1,\mathbf{Q})\) 上的投影平行于中心 \(\mathbf{Q}.1\)（\(\mathfrak{gl}(l+1,\mathbf{Q})\) 的中心），所得的是 \((\mathfrak{g},\mathfrak{h})\) 中由许可格 \(\mathrm{P}(\mathrm{R}^{\vee})\)（位于 \(\mathfrak{h}\) 中）和谢瓦莱系 \((X_{\alpha})\) 定义的谢瓦莱序。我们指出，\(\mathfrak{gl}(l+1,\mathbf{Z})\) 不是它与 \(\mathfrak{sl}(l+1,\mathbf{Q})\) 的交及其与 \(\mathfrak{gl}(l+1,\mathbf{Q})\) 的中心的交的直和。